\documentclass[12pt,a4paper]{article}

\usepackage[margin=1in]{geometry}
\usepackage{hyperref}
\usepackage{xcolor}
\usepackage{graphicx}
\usepackage{titlesec}


\pagenumbering{arabic}


\hypersetup{
colorlinks=true,
urlcolor=blue
}


\titleformat{\section}
{\large\bfseries\color{blue}}
{}
{0em}
{}[\titlerule]



\begin{document}



\begin{center}

{\Huge \textbf{Computer Science Portfolio Report}}

\vspace{10pt}

{\large B201 Computer Science Lab}

\vspace{10pt}

\textbf{Student Name: Simran}

\end{center}




\section{Introduction}


The objective of this project is to design and develop a professional computer science portfolio website that demonstrates my technical skills, academic background, and software development projects.

The portfolio is created as a digital profile for a working student application. It allows potential employers to view my education, programming knowledge, technical abilities, and practical projects in an organised and user-friendly format.

The project focuses on creating a clean, responsive, and accessible website using frontend technologies. Along with the website, a LaTeX-based CV has been created and included in the GitHub repository as required.



\section{Portfolio Links}


The completed portfolio website and source code are available through the following links:


\textbf{Portfolio Website:}

\url{https://yourusername.github.io/computer-science-portfolio/}


\vspace{8pt}


\textbf{GitHub Repository:}

\url{https://github.com/yourusername/computer-science-portfolio}




\section{Portfolio Website Structure}


The portfolio website is developed using HTML5 and CSS3.

The website contains the following sections:



\subsection{Home Section}


The home section introduces my profile as a Computer Science student and developer. It contains a short introduction along with buttons for viewing projects and downloading the CV.



\subsection{About Me Section}


This section provides information about my interests in software development, technology, and continuous learning. It explains my motivation towards improving technical skills and building practical solutions.



\subsection{Education Section}


The education section presents my academic background including:


\begin{itemize}

\item Secondary Education

\item Upper Secondary Education (Non-Medical - PCM)

\item Bachelor's Degree in Computer Science

\end{itemize}


This section helps recruiters understand my educational foundation and technical journey.



\subsection{Skills Section}


The skills section is divided into two categories:



\textbf{Technical Skills}

\begin{itemize}

\item HTML5

\item CSS3

\item Python

\item SQL

\item MySQL

\item Git and GitHub

\item VS Code

\end{itemize}



\textbf{Soft Skills}

\begin{itemize}

\item Problem Solving

\item Teamwork

\item Communication

\item Adaptability

\item Creative Thinking

\end{itemize}





\subsection{Projects Section}


The projects section showcases practical work completed using different technologies.

The included projects are:



\textbf{Guess Game}

A Python-based guessing game developed using programming logic and conditional statements. This project demonstrates basic programming concepts and problem-solving ability.



\textbf{Employee Management System}

A database project created using SQL to store, manage, and retrieve employee records efficiently.



\textbf{Frontend Website}

A responsive website developed using HTML and CSS. It demonstrates webpage structure, styling, layout design, and frontend development skills.





\section{Implementation Details}


The website implementation is divided into two main parts:



\subsection{HTML Implementation}


HTML5 is used to create the structure of the portfolio website.

Different semantic sections such as navigation, header, education, skills, projects, and contact areas were created to improve readability and maintainability.



\subsection{CSS Implementation}


CSS3 is used for styling and improving the visual appearance of the website.

The design includes:


\begin{itemize}

\item Responsive layouts

\item Card-based design

\item Navigation styling

\item Button designs

\item Proper spacing and alignment

\item Modern user interface elements

\end{itemize}



\subsection{CV Integration}


A LaTeX-created CV is integrated into the website using a download button.

Users can directly access the PDF version of the CV from the homepage.



\section{Technologies Used}


\begin{itemize}


\item \textbf{HTML5:}
Used for creating the structure and content of the portfolio.


\item \textbf{CSS3:}
Used for styling, layout design, and responsive behaviour.


\item \textbf{LaTeX:}
Used to create a professional CV and project documentation.


\item \textbf{GitHub:}
Used for version control and storing project files.


\item \textbf{GitHub Pages:}
Used for hosting the portfolio website online.


\end{itemize}





\section{Design Decisions}


The design of the portfolio was created with simplicity, readability, and professionalism in mind.

A card-based layout was selected because it presents information in separate sections, making it easier for visitors to understand education, skills, and projects.

A navigation bar was added to improve accessibility and allow users to quickly move between sections.

The website uses responsive design principles to ensure compatibility with different screen sizes including desktops and mobile devices.

The colour scheme and spacing were selected to create a clean and modern appearance suitable for a professional portfolio.



\section{GitHub Repository Structure}


The repository is organised as follows:


\begin{verbatim}

Computer-Science-Portfolio

|
|-- index.html
|
|-- style.css
|
|-- README.md
|
|-- CV
|     |-- resume.tex
|     |-- resume.pdf
|
|-- Report
      |-- portfolio_report.tex

\end{verbatim}



\textbf{index.html}

Contains the complete website structure and content.



\textbf{style.css}

Contains all styling rules including layout, colours, spacing, and responsive design.



\textbf{CV Folder}

Contains the LaTeX source file and generated PDF CV.



\textbf{Report Folder}

Contains the LaTeX project report documentation.





\section{Testing and Evaluation}


The portfolio was tested by checking:


\begin{itemize}

\item All navigation links correctly move to sections.

\item CV download button opens the PDF file.

\item Website layout works properly on different screen sizes.

\item Project and education information is displayed correctly.

\item Repository files are organised properly.

\end{itemize}



The final result is a functional and professional portfolio website suitable for presenting technical skills to employers.



\section{Strengths}


The strengths of the portfolio include:


\begin{itemize}


\item Clean and professional user interface.

\item Easy navigation between sections.

\item Responsive website design.

\item Proper organisation of GitHub repository.

\item Integration of LaTeX-based CV.

\item Clear presentation of technical projects.


\end{itemize}





\section{Limitations}


Although the portfolio meets its purpose, some limitations exist:


\begin{itemize}


\item The website does not contain backend functionality.

\item Contact section does not include a working form.

\item Projects are currently displayed without live demonstrations.

\item Limited interactive features.


\end{itemize}





\section{Future Improvements}


Possible future improvements include:


\begin{itemize}


\item Developing a React.js version of the portfolio.

\item Adding animations and advanced UI effects.

\item Creating a backend-powered contact form.

\item Adding project filtering functionality.

\item Including more advanced software projects.

\item Adding a blog section for technical articles.


\end{itemize}





\section{Conclusion}


This project successfully demonstrates the development of a computer science portfolio website using HTML5 and CSS3.

The portfolio provides a structured presentation of my education, skills, and projects while following professional design principles.

The project improved my understanding of frontend development, documentation, GitHub repository management, and professional presentation of technical work.



\section{References}


Mozilla Developer Network (2026) HTML Documentation.

Available at:
\url{https://developer.mozilla.org/}



Mozilla Developer Network (2026) CSS Documentation.

Available at:
\url{https://developer.mozilla.org/}



GitHub (2026) GitHub Pages Documentation.

Available at:
\url{https://docs.github.com/}



\end{document}